% ═══════════════════════════════════════════════════════════════════════
\section{La droite qui représente le mieux un nuage}
% ═══════════════════════════════════════════════════════════════════════

\begin{frame}{Deux façons de mesurer l'écart à une droite}
  \small
  \begin{itemize}
    \item Droite de régression des moindres carrés : minimise la somme des
      carrés des écarts verticaux. Elle répond à une question de
      prédiction : connaissant $x$, quel $y$ prévoir ? Pente
      $a=s_{xy}/s_x^2=\num{0.8}$, d'où $y=\num{0.8}\,x+19$.
    \item Prédire $x$ à partir de $y$ : écarts horizontaux, autre droite,
      $y=\num{1.25}\,x-\num{12.5}$ !
    \item Pour décrire la forme du nuage, sans privilégier de variable :
      distances perpendiculaires à la droite.
  \end{itemize}
  \FigLivre[0.48\textheight]{fig-acp-residus}
  \centering\scriptsize Figure 5.2 --- Écarts verticaux (régression) et
  distances perpendiculaires : deux droites différentes.
\end{frame}

\begin{frame}{Définition 5.3 --- Droite d'ajustement orthogonal}
  \begin{Enonce}{Définition 5.3 --- Droite d'ajustement orthogonal}
    Soit $\vec x_1,\dots,\vec x_n$ des points du plan (ou de $\R^p$), et
    $\mathcal D$ une droite. Notons $d_i$ la distance de $\vec x_i$ à sa
    projection orthogonale sur $\mathcal D$. La droite qui rend la somme
    \[
      S(\mathcal D)=d_1^2+d_2^2+\dots+d_n^2
    \]
    la plus petite possible s'appelle la \defi{droite d'ajustement
    orthogonal} du nuage, ou \defi{axe principal}.
  \end{Enonce}
  \begin{itemize}
    \item Comme les droites de régression, elle passe par le point moyen
      (admis) : sinon, on pourrait la translater pour la rapprocher des
      points.
    \item On centre donc les données, par exemple $\vec x_A=(-10,-3)$, et on
      cherche parmi les droites $\Span\{\vec u\}$, $\norme{\vec u}=1$.
  \end{itemize}
\end{frame}

\begin{frame}{Projeter sur une droite : le théorème de Pythagore}
  \begin{columns}[c]
    \begin{column}{0.55\textwidth}
      \small Projection orthogonale (définition D.4) :
      \[
        \proj_{\vec u}\vec x_i=(\vec x_i\cdot\vec u)\,\vec u=t_i\,\vec u,
        \qquad t_i=\vec x_i\cdot\vec u .
      \]
      $t_i$ : la position, le long de la droite, de l'« ombre » de
      $\vec x_i$. Triangle rectangle :
      \[
        \norme{\vec x_i}^2=t_i^2+d_i^2. \tag{5.1}
      \]
      L'hypoténuse $\norme{\vec x_i}$ (distance au point moyen) ne dépend pas
      de la droite : quand la droite pivote, si $t_i$ s'allonge, $d_i$
      raccourcit.
    \end{column}
    \begin{column}{0.43\textwidth}
      \FigLivre[0.6\textheight]{fig-acp-pythagore}
      \scriptsize Figure 5.3 --- La projection de $\vec x_i$ sur la droite
      engendrée par $\vec u$.
    \end{column}
  \end{columns}
\end{frame}

\begin{frame}{L'observation clé du chapitre}
  En additionnant (5.1) sur tous les points :
  \[
    \underbrace{\sum_{i=1}^n\norme{\vec x_i}^2}_{\text{ne dépend pas de }\vec u}
    =\sum_{i=1}^n t_i^2+\sum_{i=1}^n d_i^2 . \tag{5.2}
  \]
  Les $t_i$ ont une moyenne nulle
  ($\frac1n\sum t_i=\big(\frac1n\sum\vec x_i\big)\cdot\vec u=0$) :
  $\frac1n\sum t_i^2$ est la variance des projections.
  \begin{Enonce}{Proposition 5.1 --- Deux problèmes équivalents}
    Pour un nuage centré, la droite $\Span\{\vec u\}$ qui minimise la somme
    des carrés des distances des points à la droite est aussi celle qui
    maximise la variance des projections des points sur la droite.
  \end{Enonce}
  \small La droite passe au plus près des points $\iff$ les points
  projetés sont le plus étalés possible : l'ombre du nuage est la plus
  longue.
\end{frame}

\begin{frame}{Exemple 5.1 --- Trois directions pour les six élèves}
  \small Trois droites par le point moyen : $\vec e_1=(1,0)$,
  $\vec u=\frac1{\sqrt2}(1,1)$, $\vec w=\frac1{\sqrt2}(1,-1)$.
  \begin{itemize}
    \item $\vec e_1$ (horizontale) : $t_i=x_i-\bar x$, variance des
      projections $s_x^2=50$.
    \item $\vec u$ : $t_i=\frac1{\sqrt2}\big((x_i-\bar x)+(y_i-\bar y)\big)$,
      \[
        \sqrt2\,(t_A,\dots,t_F)=(-13,-16,-11,14,13,13),\qquad
        \frac16\sum t_i^2=\frac{1080}{12}=90 .
      \]
    \item $\vec w$ : $\sqrt2\,(t_A,\dots,t_F)=(-7,2,5,-4,-1,5)$, variance
      $\frac{120}{12}=10$.
  \end{itemize}
  \vspace{2pt}
  Variance totale : $\frac16\sum\norme{\vec x_i}^2=s_x^2+s_y^2=100$. Divisée
  par $n$, la relation (5.2) donne :
  \begin{itemize}
    \item pour $\vec u$ : $100=90+10$ (projections $90$, distances $10$) ;
    \item pour $\vec w$ : $100=10+90$.
  \end{itemize}
  La direction $\vec u$ est de loin la meilleure des trois.
\end{frame}

\begin{frame}{La variance dans une direction quelconque}
  Avec $\vec x_i=(a_i,b_i)$ et $\vec u=(u_1,u_2)$, $t_i=a_iu_1+b_iu_2$ :
  \[
    \frac1n\sum_{i=1}^n t_i^2
    =s_x^2\,u_1^2+2s_{xy}\,u_1u_2+s_y^2\,u_2^2
    =\vec u^{\mathsf T}C\,\vec u, \tag{5.3}
  \]
  où
  \[
    C=\mat{s_x^2}{s_{xy}}{s_{xy}}{s_y^2}
    \qquad\text{(six élèves : }C=\mat{50}{40}{40}{50}\text{)}.
  \]
  \begin{itemize}
    \item $C$ est la \defi{matrice de covariance} : les variances sur la
      diagonale, la covariance ailleurs.
    \item $C$ est symétrique.
  \end{itemize}
\end{frame}

\begin{frame}{Une fonction de l'angle}
  \begin{columns}[c]
    \begin{column}{0.4\textwidth}
      \small Avec $\vec u=(\cos\theta,\sin\theta)$ :
      \begin{align*}
        f(\theta)&=50\cos^2\theta+80\cos\theta\sin\theta+50\sin^2\theta\\
                 &=50+40\sin(2\theta).
      \end{align*}
      \begin{itemize}
        \item Maximum $90$ pour $\theta=45^\circ$.
        \item Minimum $10$ pour $\theta=135^\circ$.
      \end{itemize}
      La droite d'ajustement orthogonal : pente $1$, par $(70,75)$, soit
      $y=x+5$.
    \end{column}
    \begin{column}{0.58\textwidth}
      \FigLivre[0.6\textheight]{fig-acp-angle}
      \scriptsize Figure 5.4 --- Projections sur les droites d'angle
      $45^\circ$ (étalées) et $135^\circ$ (tassées) ; la variance $f(\theta)$.
    \end{column}
  \end{columns}
\end{frame}

\begin{frame}{Pour la classe 5.2 --- Une fonction sinusoïdale}
  \begin{Classe}{Une fonction sinusoïdale}
    \begin{itemize}
      \item $f(\theta)=50+40\sin(2\theta)$ : une fonction sinusoïdale de
        cinquième secondaire (SN), de valeur moyenne $50$, d'amplitude $40$
        et de période $180^\circ$.
      \item La période s'explique sans calcul : un demi-tour ramène une
        droite sur elle-même.
      \item Le problème d'optimisation est accessible à un élève du
        secondaire, pourvu qu'on lui fournisse l'expression de $f$.
      \item Activité : chercher la meilleure droite par essais, en la faisant
        tourner autour du point moyen dans GeoGebra.
    \end{itemize}
  \end{Classe}
\end{frame}

\begin{frame}{Le lien avec les vecteurs propres}
  \small $C=\mat{50}{40}{40}{50}$ a pour valeurs propres $\lambda_1=90$ et
  $\lambda_2=10$, de vecteurs propres $(1,1)$ et $(1,-1)$ (comme
  l'exemple D.7 : $C=10\mat{5}{4}{4}{5}$). L'axe principal est porté par le
  vecteur propre de la plus grande valeur propre, qui est la variance
  maximale.
  \vspace{4pt}

  Ce n'est pas une coïncidence. Pour $C$ symétrique quelconque et
  $f(\theta)=\vec u^{\mathsf T}C\,\vec u$, avec
  $\vec u'(\theta)=(-\sin\theta,\cos\theta)\perp\vec u$ :
  \[
    f'(\theta)=2\,\vec u'\cdot C\vec u .
  \]
  \begin{AuTableau}[Calcul au tableau]
    \footnotesize Dérivée d'un produit, puis symétrie :
    $\vec u^{\mathsf T}C\vec u'=\vec u'^{\mathsf T}C^{\mathsf T}\vec u=\vec u'^{\mathsf T}C\vec u$.
  \end{AuTableau}
  \begin{itemize}
    \item $f'=0\iff C\vec u\perp\vec u'\iff C\vec u$ parallèle à $\vec u$
      $\iff C\vec u=\lambda\vec u$.
    \item Dans une direction propre, $f=\vec u\cdot\lambda\vec u=\lambda$ :
      la variance est la valeur propre.
  \end{itemize}
\end{frame}

\begin{frame}{Théorème 5.2 --- La droite d'ajustement orthogonal}
  \begin{Enonce}{Théorème 5.2 --- La droite d'ajustement orthogonal}
    Soit un nuage du plan, de point moyen $M$ et de matrice de covariance
    $C$, et $\lambda_1\geq\lambda_2$ les valeurs propres de $C$, de vecteurs
    propres unitaires $\vec v_1$ et $\vec v_2$.
    \begin{enumerate}
      \item La droite d'ajustement orthogonal passe par $M$ et est dirigée
        par $\vec v_1$.
      \item La variance des projections sur cette droite vaut $\lambda_1$ ;
        la moyenne des carrés des distances à la droite vaut $\lambda_2$.
      \item La variance totale $s_x^2+s_y^2$ est égale à
        $\lambda_1+\lambda_2$.
    \end{enumerate}
  \end{Enonce}
  \begin{AuTableau}
    \footnotesize (1) Proposition 5.1 et étude de $f$ : continue et
    périodique, elle atteint son maximum là où $f'=0$, dans une direction
    propre. (3) La trace. (2) Diviser (5.2) par $n$ :
    $(s_x^2+s_y^2)-\lambda_1=\lambda_2$.
  \end{AuTableau}
\end{frame}

\begin{frame}{La forme du nuage}
  \begin{columns}[c]
    \begin{column}{0.42\textwidth}
      \small
      \begin{itemize}
        \item Variance totale $100=90+10$ : l'axe principal
          \defi{explique} $\SI{90}{\percent}$ de la variance.
        \item Remplacer chaque point par sa projection ne fait perdre que
          $\SI{10}{\percent}$ de la dispersion.
        \item Le nuage ressemble à une ellipse dirigée par les vecteurs
          propres, de demi-axes proportionnels à $\sqrt{\lambda_1}$ et
          $\sqrt{\lambda_2}$.
        \item L'axe principal $y=x+5$ est entre les deux droites de
          régression.
      \end{itemize}
    \end{column}
    \begin{column}{0.56\textwidth}
      \FigLivre[0.58\textheight]{fig-acp-axes-principaux}
      \scriptsize Figure 5.5 --- L'ellipse de demi-axes $2\sqrt{90}$ et
      $2\sqrt{10}$ ; l'axe principal et les deux droites de régression.
    \end{column}
  \end{columns}
\end{frame}

\begin{frame}{Niveau général et préférence}
  Les coordonnées sur les deux axes principaux :
  \[
    t_i=\tfrac1{\sqrt2}\big((x_i-70)+(y_i-75)\big),\qquad
    t_i'=\tfrac1{\sqrt2}\big((x_i-70)-(y_i-75)\big).
  \]
  \begin{itemize}
    \item $t_i$ est proportionnel à la somme des écarts : le niveau
      général de l'élève.
    \item $t_i'$ est proportionnel à la différence : positif pour un élève
      relativement plus fort en mathématiques qu'en sciences.
    \item Un simple changement de base orthonormée (une rotation de
      $45^\circ$ du repère) : aucune information perdue, mais
      l'information est rangée par ordre d'importance.
  \end{itemize}
\end{frame}

\begin{frame}{Remarque 5.3 --- Régression ou axe principal ?}
  \begin{Rem}{Remarque 5.3}
    \begin{itemize}
      \item L'axe principal est toujours entre les deux droites de
        régression ; les trois sont confondues seulement si les points sont
        alignés.
      \item Les droites de régression sont « tirées » vers l'horizontale ou
        la verticale : la régression vers la moyenne de Galton (taille des
        enfants et des parents).
      \item Pour prédire une variable à partir d'une autre : la régression.
        Pour décrire un nuage sans privilégier de variable : l'axe
        principal.
      \item Une faiblesse de l'axe principal : il dépend des unités de mesure
        (exercice 5.4), contrairement à la droite de régression et au
        coefficient de corrélation.
    \end{itemize}
  \end{Rem}
\end{frame}
